\documentclass[11pt]{article}
\usepackage[utf8]{inputenc}
\usepackage[T1]{fontenc}
\usepackage{graphicx}
\usepackage{amsmath}
\usepackage{amssymb}
\author{Hesam Pakdaman}
\date{\today}
\title{}

\newcommand{\collection}[2]{\{{#1}_{#2}\}}

\begin{document}

\noindent\textbf{12.} Let $\collection{G}{\alpha}$ be any open cover
of $K$. Then we know that there is an index $\alpha_0$ such that
$0 \in G_{\alpha_0}$. Since $G_{\alpha_0}$ is an open set we know that
there exists a neighborhood $N_r(0)$ with radius $r > 0$ such that
$N \subset G_{\alpha_0}$.

By archimedean property of $R$ we can find positive integers $n$ such
that $r \leq 1/n$. Since $0$ is the only limit point to $E$ and
$r > 0$ there exists a first integer\footnote{We can show that there
  exists such a first integer. If $1 > r > 0$ then $n > 1/r$ by the
  archimedean property. Simply decrement $n$ until the inequality does
  not hold. The previous number in this process is then the first
  integer. Similarly can be done with $r>1$.} $m$ such that
$r \leq 1/m$. It follows that if $n > m$ then $r > 1/n$ so that
$1/n \in G_{\alpha_0}$. This shows that there are $m$ points in $E$
that are not in $N_r$


$$\frac{1}{m}, \frac{1}{m-1}, \dots, \frac{1}{2}, 1.$$


\noindent Let $G_{\alpha_k}$ denote an open set in the collection
$\collection{G}{\alpha}$ such that $1/k \in G_{\alpha_k}$
$(k = 1, 2, 3, \dots)$.  Each of the $m$ points above can be
associated this way to at least one index in the collection
$\collection{G}{\alpha}$ because it is an open cover of $K$ and
therefore


$$ K \subset G_{\alpha_m} \cup \cdots \cup G_{\alpha_0}.$$


\noindent We have shown that any open cover of $K$ has a finite sub cover which
implies that $K$ is compact as desired.

\hfill$\square$

\end{document}